\documentclass[12pt]{article}
\usepackage[utf8]{inputenc}
\usepackage[T1]{fontenc}
\usepackage[margin=1in]{geometry}
\usepackage{setspace}
\usepackage{xcolor}
\usepackage[colorlinks=true, urlcolor=blue!60!black, linkcolor=blue!60!black]{hyperref}
\setlength{\parindent}{2em}
\setlength{\parskip}{1.25em}
\renewcommand{\baselinestretch}{1.0}
\begin{document}
\begin{center}
{\LARGE \textit{Rust v. Sullivan}, 500 U.S. 173 (1991)\par}
\vspace{0.8em}
{\large Case Brief\par}
\end{center}
\vspace{1em}

\section{Memory Jogger}

When the government appropriates public funds to establish a program promoting a particular policy, it may define the program's limits by prohibiting federally funded projects from engaging in abortion counseling, referral, or advocacy without violating the First or Fifth Amendments (\href{https://www.courtlistener.com/opinion/112593/rust-v-sullivan/}{CourtListener}; \href{https://supreme.justia.com/cases/federal/us/500/173/}{Justia}).

\section{Detailed Case Facts}

In 1970, Congress enacted Title X of the Public Health Service Act, 42 U.S.C. \S\S~300 to 300a--6, authorizing the Secretary of Health and Human Services (HHS) to make grants and enter into contracts with public or nonprofit private entities to operate voluntary family planning projects. Section 1008 of the Act, 42 U.S.C. \S~300a--6, expressly provided that ``[n]one of the funds appropriated under this subchapter shall be used in programs where abortion is a method of family planning'' (\href{https://supreme.justia.com/cases/federal/us/500/173/}{Justia}; \href{https://www.law.cornell.edu/supct/html/89-1391.ZS.html}{Cornell LII}). For roughly 18 years after enactment, HHS interpreted \S~1008 to prohibit Title X projects from performing abortions while permitting---and in later guidelines requiring---nondirective counseling and referral regarding pregnancy options, including abortion (\href{https://supreme.justia.com/cases/federal/us/500/173/}{Justia}).

In 1988, citing critical reports from the General Accounting Office and the Office of the Inspector General, the Secretary of HHS promulgated new regulations (42 C.F.R. \S\S~59.2, 59.8, 59.9, 59.10 (1989)) limiting Title X projects to preconceptional preventive family planning services and attaching three principal conditions on federal Title X funds. First, a Title X project could not ``provide counseling concerning the use of abortion as a method of family planning or provide referral for abortion as a method of family planning'' (42 C.F.R. \S~59.8(a)(1)), even upon a pregnant client's specific request; instead, projects had to refer pregnant clients for prenatal and social services promoting the welfare of mother and unborn child, and one permissible response to an abortion inquiry was that ``the project does not consider abortion an appropriate method of family planning, and therefore does not counsel or refer for abortion'' (\S~59.8(b)(5)). Second, a Title X project could not ``encourage, promote or advocate abortion as a method of family planning,'' including through lobbying, disseminating materials, providing speakers, litigating, or paying dues to groups substantially engaged in abortion advocacy (\S~59.10(a)). Third, under the ``program integrity'' rule (\S~59.9), Title X projects had to be ``physically and financially separate'' from prohibited abortion activities, maintaining objective integrity and independence through separate facilities, personnel, and accounting records rather than mere bookkeeping separation (\href{https://supreme.justia.com/cases/federal/us/500/173/}{Justia}; \href{https://www.law.cornell.edu/constitution-conan/amendment-1/selective-funding-arrangements}{Constitution Annotated}).

\section{Procedural History}

Before the 1988 regulations took effect, Title X grantees and supervising physicians---including Dr. Irving Rust, suing on behalf of themselves and their patients---along with the State and City of New York, filed two consolidated actions in the U.S. District Court for the Southern District of New York challenging the regulations on their face as unauthorized by Title X and violative of the First and Fifth Amendments. After initially granting a preliminary injunction, the District Court upheld the regulations and granted summary judgment to the Secretary (\textit{New York v. Bowen}, 690 F. Supp. 1261 (S.D.N.Y. 1988)) (\href{https://supreme.justia.com/cases/federal/us/500/173/}{Justia}).

A divided panel of the U.S. Court of Appeals for the Second Circuit affirmed (\textit{New York v. Sullivan}, 889 F.2d 401 (2d Cir. 1989)), holding that the regulations were a permissible agency construction of \S~1008 under \textit{Chevron U.S.A. Inc. v. Natural Resources Defense Council, Inc.}, 467 U.S. 837 (1984), and complied with the First and Fifth Amendments; Judge Kearse dissented in part. Meanwhile, both the First Circuit en banc (\textit{Massachusetts v. Secretary of Health and Human Services}, 899 F.2d 53 (1st Cir. 1990)) and the Tenth Circuit (\textit{Planned Parenthood Federation of America v. Sullivan}, 913 F.2d 1492 (10th Cir. 1990)) invalidated the regulations primarily on constitutional grounds. The Supreme Court granted certiorari to resolve the circuit split, heard oral argument on October 30, 1990, and affirmed the Second Circuit on May 23, 1991 (\href{https://supreme.justia.com/cases/federal/us/500/173/}{Justia}; \href{https://tile.loc.gov/storage-services/service/ll/usrep/usrep500/usrep500173/usrep500173.pdf}{U.S. Reports PDF}).

\section{Judicial Votes}

The Supreme Court affirmed by a \textbf{5--4} vote, with no concurring opinions and three separate dissents (\href{https://www.courtlistener.com/opinion/112593/rust-v-sullivan/}{CourtListener}; \href{https://supreme.justia.com/cases/federal/us/500/173/}{Justia}):

\begin{itemize}
\item \textbf{Majority (5):} Chief Justice William H. Rehnquist (author), joined by Justices Byron R. White, Anthony M. Kennedy, Antonin Scalia, and David H. Souter.
\item \textbf{Concurrences (0):} None.
\item \textbf{Dissents (4):}
  \begin{itemize}
  \item Justice Harry A. Blackmun (author), joined in full by Justice Thurgood Marshall; joined as to Part I (statutory construction and constitutional avoidance) by Justice Sandra Day O'Connor; and joined as to Parts II and III (First and Fifth Amendments) by Justice John Paul Stevens.
  \item Justice John Paul Stevens (author of separate dissent on statutory text and agency departure, while joining Parts II and III of Justice Blackmun's dissent).
  \item Justice Sandra Day O'Connor (author of separate dissent on constitutional avoidance, joining only Part I of Justice Blackmun's dissent).
  \end{itemize}
\end{itemize}

\section{Holding}

\textbf{Judgment.} The judgment of the U.S. Court of Appeals for the Second Circuit upholding the 1988 HHS Title X regulations against facial statutory and constitutional challenges was affirmed (\href{https://supreme.justia.com/cases/federal/us/500/173/}{Justia}).

\textbf{Rule Announced.} First, because \S~1008 and its legislative history are ambiguous regarding abortion counseling, referral, advocacy, and program integrity, the Secretary's regulations represent a permissible statutory construction entitled to \textit{Chevron} deference and do not raise ``grave and doubtful constitutional questions'' requiring constitutional avoidance. Second, under the First Amendment, the government may selectively fund a program to encourage certain activities (such as preconceptional family planning) and define the program's limits to exclude speech outside its scope without engaging in unconstitutional viewpoint discrimination or imposing an unconstitutional condition, so long as the condition governs the funded \textit{project} rather than barring the \textit{grantee} or its employees from speaking outside the program. Third, under the Fifth Amendment Due Process Clause, refusing to subsidize abortion counseling or referral places no governmental obstacle in a pregnant woman's path beyond her pre-existing indigency (\href{https://supreme.justia.com/cases/federal/us/500/173/}{Justia}; \href{https://www.law.cornell.edu/constitution-conan/amendment-1/selective-funding-arrangements}{Constitution Annotated}).

\textbf{Holding vs. Dicta.} The holdings are that (1) the 1988 Title X regulations are a permissible interpretation of \S~1008 and (2) they are facially valid under the First and Fifth Amendments. Important dicta appear in Chief Justice Rehnquist's caveat that government funding is not ``invariably sufficient'' to justify content control---specifically flagging traditional public forums (\textit{Kokinda}, \textit{Hague}) and universities as ``traditional sphere[s] of free expression'' (\textit{Keyishian}), and leaving open whether a comprehensive doctor-patient relationship might enjoy First Amendment protection even when government-subsidized. Likewise, the Court's later doctrinal recharacterization of \textit{Rust} as a ``government speech'' case (beginning in \textit{Rosenberger v. Rector \& Visitors of the University of Virginia}, 515 U.S. 819 (1995), and \textit{Legal Services Corp. v. Velazquez}, 531 U.S. 533 (2001)) went beyond the exact terminology of \textit{Rust} itself, which framed the case around program-scope funding limits (\href{https://supreme.justia.com/cases/federal/us/500/173/}{Justia}; \href{https://knightcolumbia.org/content/the-government-speech-doctrine-goes-to-school}{Knight First Amendment Institute}).

\section{Analysis of Opinions}

\subsection{Majority Opinion (Chief Justice Rehnquist)}

Writing for the Court, Chief Justice Rehnquist first underscored the heavy burden of a facial challenge under \textit{United States v. Salerno}, 481 U.S. 739 (1987). On the statutory issue, he concluded that \S~1008's text and legislative history were ambiguous on counseling, referral, advocacy, and physical separation, triggering deference to the Secretary's plausible construction under \textit{Chevron}. Reversing 18 years of prior agency policy did not forfeit deference because agencies possess latitude to adapt rules when supported by ``reasoned analysis,'' here supplied by GAO/OIG reports, client experience, and shifting public attitudes. Rejecting petitioners' invocation of the constitutional-avoidance canon (\textit{DeBartolo}, \textit{Catholic Bishop}), Rehnquist concluded that the regulations did not present ``grave and doubtful constitutional questions'' (\href{https://supreme.justia.com/cases/federal/us/500/173/}{Justia}).

On the First Amendment, Rehnquist relied on \textit{Maher v. Roe}, 432 U.S. 464 (1977), \textit{Harris v. McRae}, 448 U.S. 297 (1980), and \textit{Regan v. Taxation With Representation of Washington}, 461 U.S. 540 (1983), to hold that the government may make a value judgment favoring childbirth over abortion and implement it through selective funding. Refusing to fund speech outside a federal program's scope is not viewpoint discrimination or suppression of ``dangerous ideas,'' just as establishing the National Endowment for Democracy does not require funding Communism or Fascism. Distinguishing \textit{FCC v. League of Women Voters of California}, 468 U.S. 364 (1984), Rehnquist drew a sharp line between conditions placed on a \textit{grantee} (which impermissibly leverage federal money to regulate private speech outside the program) and conditions placed on a \textit{funded project}: Title X grantees and employees remained free to perform abortions and advocate for abortion through separate, non-Title X programs and in their private capacities. He also read the regulations on their face to permit emergency referrals when a woman's life is in imminent peril, and found no significant impingement on the doctor-patient relationship because Title X is limited to preconceptional care and doctors can tell patients that abortion advice is outside the program's scope. Finally, under the Fifth Amendment, Rehnquist held that \textit{McRae} and \textit{Webster v. Reproductive Health Services}, 492 U.S. 490 (1989), foreclosed any affirmative right to subsidized abortion information, distinguishing \textit{Akron} and \textit{Thornburgh} as jurisdiction-wide mandates binding all physicians (\href{https://supreme.justia.com/cases/federal/us/500/173/}{Justia}).

\subsection{Dissenting Opinion (Justice Blackmun)}

Justice Blackmun (joined in full by Justice Marshall, in Part I by Justice O'Connor, and in Parts II and III by Justice Stevens) accused the majority of casting aside constitutional avoidance to uphold, for the first time, viewpoint-based suppression of speech as a condition on public funding. In Part I, he argued that because \S~1008 was concededly ambiguous and two circuits had struck down the regulations on constitutional grounds, the Court had a plain duty under \textit{DeBartolo} and \textit{Machinists v. Street} to construe \S~1008 as barring only the performance of abortions rather than reaching grave constitutional questions (\href{https://supreme.justia.com/cases/federal/us/500/173/}{Justia}).

In Part II (First Amendment), Blackmun argued that the regulations were unmistakably content- and viewpoint-based: Title X clinics could discuss every family planning option except abortion, were forbidden from advocating abortion while anti-abortion advocacy remained unregulated, and were compelled to refer pregnant women for prenatal care and recite that the project does not consider abortion appropriate. Unlike the content-neutral lobbying restriction in \textit{Regan}, the gag rule intruded directly into truthful professional dialogue between physician and patient. Rejecting the majority's ``off-the-clock'' rationale, Blackmun cited \textit{Perry v. Sindermann}, \textit{Abood}, and \textit{Pickering} to contend that restricting employee speech inside the funded workplace still requires balancing vital speaker and patient interests against narrow government goals. In Part III (Fifth Amendment), Blackmun contended that by manipulating the physician's words in a relationship of profound trust, the government affirmatively distorted indigent women's informed choices, effectively obliterating their reproductive liberty under \textit{Roe v. Wade}, \textit{Akron}, and \textit{Thornburgh} (\href{https://supreme.justia.com/cases/federal/us/500/173/}{Justia}).

\subsection{Dissenting Opinion (Justice Stevens)}

Justice Stevens wrote separately to argue that the Court erred at \textit{Chevron} Step One: reading Title X's text and declaration of purpose as a whole---which emphasized making ``comprehensive voluntary family planning services'' and ``information'' readily available---\S~1008's ``PROHIBITION OF ABORTION'' was plainly directed at \textit{conduct} (performing abortions), not the censorship of truthful speech or medical advice. He emphasized 18 years of consistent agency construction across four presidencies and argued that in a society abhorring censorship, statutory authority to regulate conduct does not implicitly delegate authority to censor speech. Because the majority reached the constitutional issues, Stevens also joined Parts II and III of Justice Blackmun's dissent (\href{https://supreme.justia.com/cases/federal/us/500/173/}{Justia}).

\subsection{Dissenting Opinion (Justice O'Connor)}

Justice O'Connor joined Part I of Justice Blackmun's dissent and wrote separately to rest her vote solely on constitutional avoidance and judicial restraint. Observing that the regulations placed content-based speech restrictions on one of the Nation's most divisive issues and therefore raised ``serious First Amendment concerns,'' she argued that the Court should invalidate the Secretary's regulations on statutory grounds without deciding whether Congress could constitutionally enact such restrictions. Ruling on statutory grounds, she emphasized, preserves legislative democracy: Congress remains free to force the constitutional question by legislating explicitly, or to do nothing (\href{https://supreme.justia.com/cases/federal/us/500/173/}{Justia}).

\section{Examples: Future Applications}

\begin{enumerate}
\item \textbf{Same-Side (Constitutional Under \textit{Rust}): Targeted Public-Health Grant Messaging.} Congress funds a federal smoking-cessation grant program administered by private clinics and stipulates that grant-funded counselors may promote nicotine replacement therapy and abstinence, but may not recommend vaping or e-cigarettes during grant-funded sessions, while allowing the parent clinic to discuss vaping in separately funded programs. Under \textit{Rust}, this condition is constitutional because the government is defining the scope of its own public-health message rather than penalizing private speech outside the program.
\item \textbf{Same-Side (Constitutional Under \textit{Rust}): Selective Arts or Democracy Subsidies.} Congress establishes a federal grant program promoting civic literacy and constitutional democracy abroad and requires that grant-funded curricula advocate democratic self-government rather than authoritarianism. Under \textit{Rust}, the government may selectively fund a program advancing a permissible policy goal without funding competing viewpoints, provided recipients remain free to express other views using non-grant funds.
\item \textbf{Opposite-Side (Unconstitutional): Leverage Over Grantees' Private Speech Outside the Program.} Congress conditions federal global-health grants on each recipient organization adopting an organization-wide policy ``explicitly opposing prostitution and sex trafficking,'' binding the NGO even when speaking with its own private funds outside the federal program. As the Supreme Court held in \href{https://supreme.justia.com/cases/federal/us/570/205/}{\textit{Agency for International Development v. Alliance for Open Society International, Inc.}, 570 U.S. 205 (2013)}, this crosses \textit{Rust}'s line by leveraging funding to regulate speech \textit{outside} the contours of the federal program and compelling ideological affirmation.
\item \textbf{Opposite-Side (Unconstitutional): Distorting an Existing Medium of Private Advocacy.} Congress funds a legal-aid corporation to represent indigent clients in welfare-benefits litigation, but prohibits funded attorneys from arguing in court that existing welfare statutes are unconstitutional or violate federal law. As the Court held in \href{https://supreme.justia.com/cases/federal/us/531/533/}{\textit{Legal Services Corp. v. Velazquez}, 531 U.S. 533 (2001)}, \textit{Rust} does not apply where the government funds private speakers (lawyers advocating for clients) rather than its own message, and where the restriction distorts the traditional functioning of the judiciary and leaves indigent litigants with no alternative channel to raise constitutional defenses (\href{https://firstamendment.mtsu.edu/article/legal-services-corp-v-velazquez/}{First Amendment Encyclopedia}).
\item \textbf{Fence-Sitter: Speech Conditions Attached to Public University Research Grants.} A federal agency awards biomedical research grants to public university professors to study climate-change adaptation, but stipulates that no funded researcher may mention fossil-fuel phaseouts in grant-funded publications or campus lectures connected to the grant. Because Chief Justice Rehnquist expressly flagged universities as ``traditional sphere[s] of free expression'' where funding conditions face heightened First Amendment limits (\textit{Keyishian}), this scenario sits squarely on the contested frontier between \textit{Rust}'s program-definition rule and academic-freedom protections against viewpoint-based censorship.
\end{enumerate}

\section{Critique}

\textbf{Originalist, Textualist, and Libertarian Perspectives.} Conservative defenders of \textit{Rust} emphasize democratic accountability and Spending Clause realism: if elected officials could not control the message of government-funded programs, the government could not govern effectively, and taxpayers would be forced to subsidize viewpoints they reject (\href{https://fedsoc.org/case/rust-v-sullivan}{Federalist Society}; \href{https://www.cardozolawreview.com/the-end-of-government-speech/}{Cardozo Law Review}). At the same time, libertarian and textualist scholars have raised serious reservations. Libertarian critics of weak unconstitutional-conditions scrutiny warn that as the modern federal spending state expands, allowing the government to buy silence through conditional grants erodes First Amendment liberty (\href{https://www.cato.org/supreme-court-review/2005-2006/unanimously-wrong}{Cato Supreme Court Review}). On the administrative-law side, textualists who later critiqued \textit{Chevron} deference (and celebrated its eventual overruling) note that \textit{Rust}'s statutory holding relied heavily on deferring to an executive agency's sharp policy reversal over Justice Stevens's textual reading and Justice O'Connor's constitutional-avoidance dissent (\href{https://firstamendment.mtsu.edu/article/rust-v-sullivan/}{First Amendment Encyclopedia}).

\textbf{Progressive and Doctrinal Perspectives.} Progressive and First Amendment scholars---including Dorothy Roberts, David Cole, and Robert Post---have forcefully critiqued \textit{Rust} for enabling ``government control of knowledge'' and distorting professional spheres of trust (\href{https://knightcolumbia.org/content/the-government-speech-doctrine-goes-to-school}{Knight First Amendment Institute}; \href{https://www.fire.org/news/harvard-law-review-note-tackles-complex-issue-distinguishing-government-versus-private-speech}{FIRE}). Critics argue that \textit{Rust}'s subsequent transformation into the ``government speech doctrine'' allows the state to co-opt private professionals (such as doctors) into delivering state-scripted ideology while concealing the government's hand from vulnerable, low-income patients who reasonably expect candid medical advice (\href{https://www.fire.org/research-learn/what-government-speech-doctrine}{FIRE}; \href{https://www.cardozolawreview.com/the-end-of-government-speech/}{Cardozo Law Review}). Others note the deep instability of the line between defining a program's scope (\textit{Rust}) and impermissibly imposing viewpoint-based conditions (\textit{Velazquez}, \textit{AOSI}).

\textbf{Underlying Competing Values.} \textit{Rust} exposes the core tension between \textit{democratic fiscal control} (the elected branches' power to decide what policies and messages public funds will promote) and \textit{individual expressive autonomy and professional integrity} (protecting speakers, doctors, and low-income recipients from state-engineered ideological distortion in the marketplace of ideas).

\section{Key Quotations}

\begin{quote}
``The Government can, without violating the Constitution, selectively fund a program to encourage certain activities it believes to be in the public interest, without at the same time funding an alternate program which seeks to deal with the problem in another way. In so doing, the Government has not discriminated on the basis of viewpoint; it has merely chosen to fund one activity to the exclusion of the other.''
\end{quote}
\noindent --- Chief Justice Rehnquist, writing for the Court, 500 U.S. at 193 (\href{https://supreme.justia.com/cases/federal/us/500/173/}{Justia}; \href{https://www.courtlistener.com/opinion/112593/rust-v-sullivan/}{CourtListener}).

\begin{quote}
``When Congress established a National Endowment for Democracy to encourage other countries to adopt democratic principles, 22 U.S.C. \S~4411(b), it was not constitutionally required to fund a program to encourage competing lines of political philosophy such as Communism and Fascism.''
\end{quote}
\noindent --- Chief Justice Rehnquist, writing for the Court, 500 U.S. at 194 (\href{https://supreme.justia.com/cases/federal/us/500/173/}{Justia}).

\begin{quote}
``The Secretary's regulations do not force the Title X grantee to give up abortion-related speech; they merely require that the grantee keep such activities separate and distinct from Title X activities.''
\end{quote}
\noindent --- Chief Justice Rehnquist, writing for the Court, 500 U.S. at 196 (\href{https://supreme.justia.com/cases/federal/us/500/173/}{Justia}).

\begin{quote}
``Until today, the Court never has upheld viewpoint-based suppression of speech simply because that suppression was a condition upon the acceptance of public funds. Whatever may be the Government's power to condition the receipt of its largess upon the relinquishment of constitutional rights, it surely does not extend to a condition that suppresses the recipient's cherished freedom of speech based solely upon the content or viewpoint of that speech.''
\end{quote}
\noindent --- Justice Blackmun, dissenting, 500 U.S. at 207 (\href{https://supreme.justia.com/cases/federal/us/500/173/}{Justia}; \href{https://www.law.cornell.edu/supct/html/89-1391.ZD1.html}{Cornell LII}).

\begin{quote}
``If we rule solely on statutory grounds, Congress retains the power to force the constitutional question by legislating more explicitly. It may instead choose to do nothing. That decision should be left to Congress; we should not tell Congress what it cannot do before it has chosen to do it.''
\end{quote}
\noindent --- Justice O'Connor, dissenting, 500 U.S. at 224 (\href{https://supreme.justia.com/cases/federal/us/500/173/}{Justia}; \href{https://tile.loc.gov/storage-services/service/ll/usrep/usrep500/usrep500173/usrep500173.pdf}{U.S. Reports PDF}).

\end{document}
